
\appendix

\section{Pattern classes}
\label{app:classes}

The 299 patterns are grouped into nine classes by what the closing message is
doing. A class name written with a trailing question mark marks a doubt, which
defers to the judge instead of deciding.

\begin{table}[h]
\centering
\begin{tabular}{lrp{4.0cm}p{4.4cm}}
\toprule
Class & Patterns & Reads as & Example trigger \\
\midrule
\texttt{announce} & 70+15 & a step named instead of taken & ``next I will'', \emph{paso a} \\
\texttt{ask} & 56+1 & a question it may answer itself & ``do you want me to'', \emph{quer\'es que} \\
\texttt{defer} & 32+2 & work made conditional on being asked & ``if you want, I can'' \\
\texttt{wait?} & 0+33 & deferring to the user's clock & ``I'll wait'', ``let me know when'' \\
\texttt{handoff} & 24 & giving the work back & ``over to you'', ``your call'' \\
\texttt{errand?} & 0+23 & sending the user to run something & ``run this and tell me'' \\
\texttt{read} & 21+2 & asking what the repo already answers & ``where is the spec'' \\
\texttt{idle} & 12+1 & reporting that nothing happened & ``still running'', ``no changes'' \\
\texttt{ack} & 7 & a bare acknowledgement & ``done'', ``ok'' \\
\bottomrule
\end{tabular}
\caption{The nine pattern classes. The second column counts firm patterns plus doubts, where a doubt raises the message to the judge instead of deciding. 222 of the 299 are firm.}
\end{table}

Two classes need an exception that is decided outside the patterns.
\texttt{ask} patterns are suppressed on long messages, because a long report
that ends in a courtesy question is a report, not a request for permission.
\texttt{wait?} patterns are suppressed when the transcript shows the agent is
waiting on work it launched itself.

\texttt{wait?} and \texttt{errand?} are written as doubts because they fire
below the base rate of the failure they look for: 0.034 and 0.050 against
0.104 (Table~\ref{tab:clean}). They raise the closing to the judge and decide
nothing on their own.

\section{The judge's instruction}
\label{app:prompt}

The judge's system prompt is a decision rule, not a persona. Its shape, in
order: the agent has standing permission and never needs to ask; a definition of
\textsc{stop} listing the observable behaviours; a definition of \textsc{ok}; a
paragraph on the most common hard case, in which the message reports real
finished work \emph{and} names something still open, which is \textsc{stop}; a
counterweight stating that a closed caveat is not a pending item, because an
agent that learns to delete its warnings to get past the gate is worse than one
that stops; a note that deferred work often dresses as a disclosure; a note that
a conditional offer is a stop; a short reading list for River Plate Spanish; and
a test in place of a default: is there a concrete action the agent could take
right now without information only the developer has, with \textsc{ok} when
unsure. That last line was measured against its two alternatives in
Table~\ref{tab:ab}.

The counterweight paragraph was added after the gate began punishing useful
warnings. It is the only part of the prompt that argues against the gate's own
purpose, and it is load-bearing.

\section{The grading rule}
\label{app:grading}

\begin{quote}
\ttfamily\small
for each logged block:\\
\hspace*{2em}find the blocked message in its transcript\\
\hspace*{2em}count tool calls until the user speaks again\\
\hspace*{2em}(a tool result is not the user speaking)\\
\hspace*{2em}fewer than 3 tool calls $\Rightarrow$ the block bought nothing\\
\\
for each pattern class:\\
\hspace*{2em}if it has 8 or more graded blocks\\
\hspace*{2em}and 70\% or more bought nothing:\\
\hspace*{4em}name it for demotion
\end{quote}

The threshold of three tool calls exists because a woken agent commonly reads
one or two files to re-establish context before deciding to stop again. Zero is
too strict and would count that re-orientation as work.

\section{Reproducing the numbers}

Every table in this paper is produced by a script in the hook's own
repository, run against the transcripts on the developer's machine. The
corpus builder is the only one that calls the model; the rest read its stored
labels and are arithmetic, so they run in seconds and give the same answer
every time.

\begin{center}
\begin{tabular}{lp{6.2cm}p{3.2cm}}
\toprule
Script & What it does & Feeds \\
\midrule
\texttt{corpus.py} & builds the labelled corpus; \texttt{--patterns} re-scores the pattern lane alone and keeps the stored verdicts & Table~\ref{tab:clean} inputs \\
\texttt{bench-lanes.py} & lane and class ablation with Wilson intervals & Table~\ref{tab:clean} \\
\texttt{bench-split.py} & the same search scored on held-out sessions & transfer result \\
\texttt{bench-prompt.py} & three closing instructions over the same rows & Table~\ref{tab:ab} \\
\texttt{bench-sure.py} & the judge's confidence against the label & confidence bands \\
\texttt{bench-contradiction.py} & delivery claims against the commands that ran & claim verification \\
\texttt{judge-report.py} & grades each block by the tool calls it bought & block conversion \\
\texttt{test\_paper.py} & fails when a number here stops matching the repo & all of it \\
\bottomrule
\end{tabular}
\end{center}

The corpus is versioned by keeping the file from each build, so a table can be
recomputed against the state it was written from. The builder takes about 13
minutes over 764 closings on a consumer GPU,
one model call per closing at temperature 0 with a fixed seed. The prompt
ablation is three passes of the same shape. Nothing else here needs a GPU.
